%% ----- Seitenränder ---------------------------------------------------------
%\setlength{\topskip}{\ht\strutbox} %% behebt Warnung von geometry
\ifthenelse{\boolean{isDebug}}{
    \geometry{showframe}
}{}
\ifthenelse{\boolean{isBookPrint}}{}{
    %% non-book print - manuall adapt to DES template
    \geometry{
        left=35mm,
        right=25mm,
        top=40mm,
        bottom=40mm,
        includehead=false, %% head as content, stay within page borders
        includefoot=false, %% foot as content, stay within page borders
    }
}

%% ----- Kopf- und Fußzeilen --------------------------------------------------
\ifthenelse{\boolean{isBookPrint}}{
    \usepackage[
        automark, %% \chaptermark und \sectionmark automatisch definiert
    ]{scrlayer-scrpage}
    \pagestyle{scrheadings}
    \clearmainofpairofpagestyles

    \lehead
    {\small\textsc{\color{gray}\hspace{-1.5em}\leftmark}}
    \rohead
    {\small\textsc{\color{gray}\hspace{-1.5em}\rightmark}}

    \ofoot
    [\large\textsc{-~\pagemark~-}]
    {\large\textsc{-~\pagemark~-}}

    \ifthenelse{\boolean{isFinalVersion}}{}{
        \cfoot
        [\large\textsc{\color{red}-~\version~-}]
        {\large\textsc{\color{red}-~\version~-}}
    }
}{
    \usepackage[
        headsepline, %% Trennlinie unter Kopfzeile
        footsepline, %% Trennlinie über Fußzeile
        plainheadsepline, %% Kopftrennlinie auf chapter seiten
        plainfootsepline, %% Fußtrennlinie auf chapter seiten
    ]{scrlayer-scrpage}
    \pagestyle{scrheadings}
    \clearpairofpagestyles

    \renewcommand{\chaptermark}[1]{\markright{\thechapter~#1}\markleft{\thechapter~#1}}
    \renewcommand{\sectionmark}[1]{\markright{\thesection~#1}}
    \addtokomafont{headsepline}{\color{gray}}
    \addtokomafont{footsepline}{\color{gray}}

    \ihead
    {\small\textsc{\color{gray}\\\leftmark}}
    \ohead
    {\small\textsc{\color{gray}\\\rightmark}}
    \ifoot
    [\small\textsc{\color{gray}\paperTitleCompact\\\student}]
    {\small\textsc{\color{gray}\paperTitleCompact\\\student}}
    \ofoot
    [\large\textsc{\\-~\pagemark~-}]
    {\large\textsc{\\-~\pagemark~-}}

    \ifthenelse{\boolean{isFinalVersion}}{}{
        \chead
        [\large\textsc{\color{red}\version\\}]
        {\large\textsc{\color{red}\version\\}}
    }
}

%\setlength{\headheight}{15mm} %% Höhe der Kopfzeile
%\setheadwidth[0pt]{textwithmarginpar} %% enlarge Head over Text (falls Logo den Text überdeckt)
%\setlength{\footskip}{10mm} %% Höhe Fußzeile: abstand inhalt zur "grundlinie" (Textboden); nicht obere Kante vom Text
%\setlength{\headsep}{50mm} %% separation space between header and text
%\setlength{\footnotesep}{50mm} %% separation space between text and footnote area


%% ----- Font Optimierung -----------------------------------------------------
%% below settings are from DES template
\usepackage[
    activate={true,nocompatibility}, %% activate protrusion and expansion
    final, %% enable microtype even if document option "draft" is used
    tracking=true,
    kerning=true,
    spacing=true,
    factor=1100, %% add 10% to the protrusion amount (default is 1000)
    stretch=10, %% reduce stretchability (default is 20)
    shrink=10, %% reduce shrinkability (default is 20)
]{microtype}
\SetExtraKerning[unit=space] {
    encoding={*},
    family={bch},
    series={*},
    size={footnotesize,small,normalsize}
} {
    \textendash={400,400}, %% en-dash, add more space around it
    "28={ ,150}, %% left bracket, add space from right
    "29={150, }, %% right bracket, add space from left
    \textquotedblleft={ ,150}, %% left quotation mark, space from right
    \textquotedblright={150, } %% right quotation mark, space from left
}


%% ----- Zeilenabstände -------------------------------------------------------
\onehalfspacing %% Zeilenabstand 1,5 Zeilen
\frenchspacing %% erzeugt ein wenig mehr Platz hinter einem Punkt


%% ----- Top-Level Kapitel ----------------------------------------------------
%% format: <kapitelNr> | <kapitelTitel>
%% wobei <kapitelNr> rießig ist
\newfont{\chapterNumber}{eurb10 scaled 4000}
\renewcommand*{\chapterformat}{%
    \chapterNumber\thechapter%
    \hspace{0.3em}%
    \smash{\protect\rule[-.15em]{1pt}{1em}}%
    \hspace{0.3em}%
}
\ifthenelse{\boolean{isBookPrint}}{}{
    \RedeclareSectionCommand[beforeskip=1em]{chapter}
}


%% ----- Font-Stil der Überschriften -----------------------------------------------
\setkomafont{sectioning}{\normalfont\normalcolor\bfseries} %% in TOC
\SetTracking[context=largetracking]{encoding=*}{160}
\SetTracking[context=smalltracking]{encoding=*}{80}
\setkomafont{chapter}{\microtypecontext{tracking=largetracking}\normalfont\LARGE\lsstyle\bfseries}
\setkomafont{section}{\microtypecontext{tracking=smalltracking}\normalfont\Large\lsstyle\bfseries}
\setkomafont{subsection}{\microtypecontext{tracking=smalltracking}\normalfont\Large\lsstyle}
\setkomafont{subsubsection}{\microtypecontext{tracking=smalltracking}\normalfont\large\lsstyle}
%% Stil der Abschnittsnummerierung %% TODO [STYLE] make section numbers bold
%\renewcommand\sectionformat{\bfseries\thesection\autodot\enskip}
%\renewcommand\subsectionformat{\bfseries\thesubsection\autodot\enskip}
%\renewcommand\subsubsectionformat{\bfseries\thesubsubsection\autodot\enskip}

%% ----- Schusterjungen und Hurenkinder vermeiden -----------------------------
%% TODO set clubpenalty and widowpenalty
%\clubpenalty=10000
%\widowpenalty=10000


%% ----- Quellcode-Ausgabe formatieren ----------------------------------------
\lstset{numbers=left, numberstyle=\tiny, numbersep=5pt, breaklines=true}
\lstset{emph={square}, emphstyle=\color{red}, emph={[2]root,base}, emphstyle={[2]\color{blue}}}

\counterwithout{footnote}{section} %% Fußnoten fortlaufend durchnummerieren
\setcounter{tocdepth}{2} %% im Inhaltsverzeichnis werden die Kapitel bis zum Level der paragraphen übernommen
\setcounter{secnumdepth}{4} %% Kapitel bis zum Level der paragraphen werden nummeriert
%% \subsubsectionlevel = level 3


%% ----- Aufzählungen anpassen ------------------------------------------------
\renewcommand{\autodot}{} %% Remove all end-of-counter dots
\renewcommand{\labelenumi}{\arabic{enumi}.}
\renewcommand{\labelenumii}{\arabic{enumi}.\arabic{enumii}.}
\renewcommand{\labelenumiii}{\arabic{enumi}.\arabic{enumii}.\arabic{enumiii}}

%% ----- Bezeichner für Kapitel über \autoref ---------------------------------
%% (original in babel ngerman: Kapitel, Abschnitt, Unterabschnitt, Unterunterabschnitt)
\addto\extrasngerman{
    \renewcommand{\sectionautorefname}{\sectionPrefix}
    \renewcommand{\subsectionautorefname}{\sectionPrefix}
    \renewcommand{\subsubsectionautorefname}{\sectionPrefix}
}

%% ----- Farben ---------------------------------------------------------------
\newcommand{\colorLinks}{true}
\newcommand{\linkColor}{linkColor}
\ifthenelse{\boolean{isBookPrint}}{
    \renewcommand{\linkColor}{linkColorBook}
}{}

%% ----- Farben ---------------------------------------------------------------
\definecolor{main}{RGB}{0, 64, 108}
\definecolor{linkColor}{RGB}{0, 64, 108}
\definecolor{linkColorBook}{RGB}{0, 37, 61}
\definecolor{evenRow}{rgb}{0.9, 0.9, 0.9}

%% ----- für Code Schemas -----
%% TODO find better way of producing color schema for code/listings
